% Dear Hiring Committee, 


% I am writing to express my interest in the
% % job title
% Postdoctoral Researcher position (reference number: OT-30574) on Hybrid Mechanistic \& AI Modeling of Microbiomes.
% It is with great enthusiasm that I submit my application for this position, as it matches to a great extent with my research trajectory at the intersection of microbial ecology, metabolic modeling, bioinformatics, and systems biology, while offering an exciting opportunity to further develop my expertise in dynamic and predictive modeling of microbial communities.

% % current position
% I am currently a postdoctoral researcher at the \href{https://m3i-lab.github.io//}{Mechanisms of microbial metabolic interactions (m3i) lab}, at the Cluster of Excellence
% \href{https://www.microverse-cluster.de/en/}{"Balance of the Microverse"}, part of the 
% \href{https://www.uni-jena.de/en}{Friedrich-Schiller University Jena}.  
% % role
% My current research focuses on resource hierarchies inference and their use in predicting microbial interactions and community dynamics.
% To this end, I am developing and using both topological and constraint-based metabolic modeling approaches, integrating multi-omics data. 
% Moreover, it serves as my first endeavor in bringing together metabolic modeling approaches with mathematical modeling, since we are now working on modeling a five-species community, aiming to decipher its dynamics considering diauxic shifts. 

% I bring a strong background in microbial ecology, bioinformatics, and metabolic modeling.
% Since my PhD, I have been contributing on the \href{https://geomscale.github.io/}{GeomScale project} on the development of random flux sampling methods for genome-scale metabolic models~\cite{socg2021,chalkis2024dingo}.
% During my first postdoctoral position, I contributed in the development of \texttt{DNNGIOR}, the deep-learning based framework for gapfilling genome-scale metabolic models~\cite{boer2024improving} leading to a particularly direct connection to this position.
% As a matter of fact, a few months after completing this project, it was the first time for me to visit INRAE in the framework of the \href{https://ai-microbiome-school.onrender.com/}{School of Artificial Intelligence Applied to Microbiomes}.
% At the same time, I also developed a \href{https://microbetag.readthedocs.io/}{microbial co-occurrence network annotator tool}~\cite{microbetag2025}; among its features, \texttt{microbetag} makes use of the \textit{seed} approach, a topological method for predicting metabolic interactions between microbial species, to annotate microbial co-occurrence networks with potential metabolic interactions. 
% Further, I deepened my understanding on the limitations of classic FBA-based approaches in interactions inference~\cite{joseph2024predicting}. 
% To support further our metabolic modeling based efforts, I have also been working on 
% the design and development of $\mu$GrowthDB~\cite{mgrowthdb_bibrxiv}, a \href{https://mgrowthdb.gbiomed.kuleuven.be/}{microbial growth data repository}; a key resource for constraining models more accurately. 

% During my PhD, I also focused on developing workflows and software tools tailored to real-world microbial communities, particularly for the analysis of high-throughput sequencing data~\cite{zafeiropoulos2020pema,zafeiropoulos2021darn,zafeiropoulos2023metagoflow}. 
% I also served as a lead developer of a \href{https://prego.hcmr.gr/}{data integration platform} designed to extract and consolidate associations among species, environments, and biological processes from both literature and abundance tables from key platforms such as MGnify, and IMG/JGI~\cite{prego2023}.

% % why i am interested in the specific position
% The position you offer strongly aligns with my research interests and would provide an excellent opportunity to continue my work on microbial interactions through metabolic modeling. 

% I am particularly interested in the project's combination of mechanistic and data-driven approaches. 
% The integration of flux balance analysis with consumer–resource models, compositional Neural ODEs, and Physics-Informed Neural Networks is appealing to me because it provides a way to combine the predictive capacity of machine learning with biological constraints and mechanistic interpretation. 
% I would be very motivated to contribute my experience in microbial systems biology and metabolic modeling while developing further expertise in dynamic modeling and physics-informed machine learning.


% Finally, I am excited by the interdisciplinary environment created by MaIAGE, PROSE, and MICALIS. The combination of expertise in mathematical biology, metabolic ecosystem modeling, microbial ecology, and fermentation provides an ideal setting for developing predictive models that can ultimately be tested and used to steer microbial communities. I would particularly value the opportunity to contribute to this effort while further developing my research profile towards quantitative and predictive microbial systems biology.



% I am confident that I can apply my current expertise to contribute effectively to the project’s goals, while also acquiring new skills in the process in a nicely put in-line with my interests and future perspectives way.
% I am enclosing my CV and a letter of recommendation from Prof. Faust. 
% Please do not hesitate to contact me for anything that might be needed. 
% Thank you for considering my application. I would be pleased to discuss how my background and research interests could contribute to the project.


% ------------
Dear Hiring Committee,

I am writing to express my interest in the Postdoctoral Researcher position (reference number: \href{https://jobs.inrae.fr/en/ot-30574}{OT-30574}) on Hybrid Mechanistic \& AI Modeling of Microbiomes. I am particularly enthusiastic about this opportunity, as it closely matches my research trajectory at the intersection of microbial ecology, metabolic modeling, bioinformatics, and systems biology, while offering an exciting opportunity to further develop my expertise in dynamic and predictive modeling of microbial communities.

I am currently a postdoctoral researcher at the \href{https://m3i-lab.github.io//}{Mechanisms of microbial metabolic interactions (m3i) lab}, at the Cluster of Excellence
\href{https://www.microverse-cluster.de/en/}{\enquote{Balance of the Microverse}}, part of the 
\href{https://www.uni-jena.de/en}{Friedrich-Schiller University Jena}.  
My current research focuses on inferring resource hierarchies, i.e. the ordered use of carbon sources by different bacteria, and using them to predict microbial interactions and community dynamics. 
To this end, I combine topological and constraint-based metabolic modeling approaches with multi-omics data.
 
Since my PhD, I have developed a broad experience in metabolic modeling, covering genome-scale model reconstruction, constraint-based analysis, flux-space sampling, and topological approaches. 
As part of the \href{https://geomscale.github.io/}{GeomScale project} project, I have worked on the development of random flux sampling methods for genome-scale metabolic models (GEMs)~\cite{socg2021,chalkis2024dingo}.

During my first postdoctoral position, I further broadened this experience across complementary aspects of metabolic modeling. First, I contributed to the development of \texttt{DNNGIOR}, a deep-learning-based framework for gap-filling GEMs~\cite{boer2024improving}, providing a particularly direct connection to the proposed work on AI-assisted metabolic model reconstruction. 
This work also gave me the opportunity to visit INRAE, in the framework of the \href{https://ai-microbiome-school.onrender.com/}{School of Artificial Intelligence Applied to Microbiomes}, where I presented my work on GEM reconstruction and flux-space sampling.

I subsequently investigated the limitations of classical FBA-based approaches for inferring microbial interactions~\cite{joseph2024predicting}, which motivated me to explore complementary approaches to metabolic interaction prediction. As part of this effort, I developed \href{https://microbetag.readthedocs.io/}{\texttt{microbetag}}, a microbial co-occurrence network annotation tool~\cite{microbetag2025}. Among its features, \texttt{microbetag} uses the \textit{seed} approach, a topological method for predicting potential metabolic interactions between microbial species, to annotate microbial co-occurrence networks with metabolic information.

Together, these experiences have given me a broad perspective on GEM-based modeling, from reconstruction and flux-space exploration to constraint-based and topological approaches for studying microbial interactions, which I see as highly relevant to the multi-scale modeling framework proposed in this position.

More recently, I have been working on the design and development of \href{https://mgrowthdb.gbiomed.kuleuven.be/}{$\mu$GrowthDB}~\cite{mgrowthdb_bibrxiv}, a microbial growth data repository intended to provide experimental information that can be used to constrain metabolic models more accurately. 
This work may prove particularly relevant to the proposed integration of high-resolution experimental data with mechanistic models and machine-learning approaches.

My background also includes extensive experience in developing computational workflows and resources for real-world microbial communities. 
During my PhD, I developed tools for the analysis of high-throughput sequencing data~\cite{zafeiropoulos2020pema,zafeiropoulos2021darn,zafeiropoulos2023metagoflow} and served as a lead developer of \href{https://prego.hcmr.gr/}{PREGO}, a data-integration platform designed to consolidate associations among microbial species, environments, and biological processes from scientific literature and large public resources such as MGnify and IMG/JGI~\cite{prego}. 
These experiences have given me a strong foundation for integrating heterogeneous biological datasets and translating them into computationally tractable representations.

This position strongly aligns with my research interests and would provide an excellent opportunity to continue developing my work on microbial interactions through metabolic and mathematical modeling. 
I am particularly attracted by the project's combination of mechanistic and data-driven approaches. 
The integration of constraint-based with consumer–resource models, compositional Neural ODEs, and Physics-Informed Neural Networks offers a compelling framework for addressing microbiome engineering challenges.
I would be highly motivated to contribute my existing expertise in microbial systems biology and metabolic modeling while developing further expertise in dynamic modeling and physics-informed machine learning.

I am confident that my experience in metabolic modeling, metabolic interaction inference, microbial ecology and multi-omics integration would allow me to contribute effectively to the project's goals. 
At the same time, the position would allow me to acquire new expertise in dynamic and AI-based modeling that fits naturally with my research interests and my long-term goal of developing predictive approaches for complex microbial ecosystems.

I am enclosing my CV. 
Please do not hesitate to contact me should you require any further information.
Thank you for considering my application. 

